% RESUMO - PT
\begin{resumo}
  \vspace{-15pt}

  O balanceamento de personagens em jogos de luta enfrenta uma tensão: equilibrar o elenco tende a homogeneizá-lo, e um elenco homogêneo perde as identidades que distinguem o gênero. Este trabalho investiga se um algoritmo genético consegue equilibrar cinco arquétipos de jogo de luta (\textit{Rushdown}, \textit{Zoner}, \textit{Grappler}, \textit{Combo Master} e \textit{Turtle}) sem destruir as suas identidades funcionais. Foi construído um simulador de combate um-contra-um em que cada atributo e cada peso comportamental afeta o desfecho das lutas, e formulada uma função de aptidão com um termo de equilíbrio e um termo de desvio em relação aos perfis canônicos, sem codificar quem deve vencer cada confronto. Um algoritmo genético escalar, o NSGA-II e um híbrido dos dois foram comparados em 20 execuções cada, com o mesmo orçamento, por testes de Wilcoxon pareados com correção de Holm. A identidade funcional foi medida por réguas comportamentais que a busca não vê, contra modelos nulos e contra um controle sem o termo de identidade. O algoritmo escalar alcança o equilíbrio em todas as execuções, com elencos tão equilibrados quanto a simetria perfeita, e o equilíbrio replica fora do laço de otimização. Sem o termo de identidade, a identidade cai ao nível de elencos sorteados, sem ganho significativo de equilíbrio; com ele, a identidade funcional fica acima do acaso, mas longe da canônica, e a política dos personagens é o componente menos preservado. O NSGA-II preserva mais identidade, mas deixa confrontos desequilibrados. O híbrido, com o mesmo orçamento, alcança a identidade do NSGA-II com o equilíbrio do algoritmo escalar, o que mostra que parte da perda de identidade observada decorre da seleção escalar sob avaliação ruidosa, e não do problema em si. Conclui-se que o equilíbrio pode ser alcançado sem destruir as identidades funcionais, e que a sua preservação depende da pressão explícita de identidade e do modo como a busca lida com o ruído.

  \textbf{Palavras-chave}: \palavraschaveemlinha
\end{resumo}


% RESUMO - EN
\begin{resumo}[Abstract]
  \vspace{-15pt}

  \begin{otherlanguage*}{english}
  Balancing characters in fighting games faces a tension: balancing a roster tends to homogenize it, and a homogeneous roster loses the identities that define the genre. This work investigates whether a genetic algorithm can balance five fighting-game archetypes (Rushdown, Zoner, Grappler, Combo Master and Turtle) without destroying their functional identities. A one-versus-one combat simulator was built in which every attribute and behavioral weight affects the outcome of fights, and a fitness function was formulated with a balance term and a term for deviation from the canonical profiles, without encoding who should win each matchup. A scalar genetic algorithm, NSGA-II and a hybrid of both were compared over 20 runs each, under the same budget, using paired Wilcoxon tests with Holm correction. Functional identity was measured by behavioral metrics that the search never sees, against null models and against a control without the identity term. The scalar algorithm reaches balance in every run, with rosters as balanced as perfect symmetry, and the balance replicates outside the optimization loop. Without the identity term, identity drops to the level of random rosters, with no significant gain in balance; with it, functional identity stays above chance but far from the canonical one, and the characters' policy is the least preserved component. NSGA-II preserves more identity but leaves unbalanced matchups. The hybrid, under the same budget, reaches the identity of NSGA-II with the balance of the scalar algorithm, showing that part of the observed identity loss stems from scalar selection under noisy evaluation rather than from the problem itself. Balance can therefore be achieved without destroying functional identities, and preserving them depends on explicit identity pressure and on how the search handles noise.

  \textbf{Keywords}: \inlinekeywords
\end{otherlanguage*}
\end{resumo}
